\ExerciseSection

\begin{enumerate}
\def\labelenumi{\arabic{enumi})}
\tightlist
\item
  Let X be a set, let Y be a uniform space containing more than one point, and let S be a non-empty set of non-empty subsets of X. Let \(Y' \subset \mathcal{F}(X; Y)\) be the set of all constant mappings of X into Y. For each \(y \in Y\) , let \(c_{y}\) be the constant mapping of X into Y whose value is y.
\end{enumerate}

\begin{enumerate}
\def\labelenumi{\alph{enumi})}
\item
  Show that \(y \to c_{y}\) is an isomorphism of Y onto the uniform subspace \(Y'\) of \(\mathcal{F}_{\mathfrak{S}}(\mathbf{X}; \mathbf{Y})\) .
\item
  \(\mathcal{F}_{\mathfrak{S}}(\mathbf{X};\mathbf{Y})\) is Hausdorff if and only if Y is Hausdorff and \(\mathfrak{S}\) is a covering of X.
\item
  \(\mathbf{Y}'\) is closed in \(\mathcal{F}_{\mathfrak{S}}(\mathbf{X};\mathbf{Y})\) if and only if \(\mathcal{F}_{\mathfrak{S}}(\mathbf{X};\mathbf{Y})\) is Hausdorff.
\end{enumerate}

\begin{enumerate}
\def\labelenumi{\arabic{enumi})}
\setcounter{enumi}{1}
\tightlist
\item
  Let X be a set, let Y be a Hausdorff uniform space containing more than one point, and let \(\mathfrak{S}_1, \mathfrak{S}_2\) be two sets of subsets of X which satisfy conditions \((\mathrm{F}_{\mathrm{I}}^{\prime}), (\mathrm{F}_{\mathrm{II}}^{\prime})\) of no. 2. Show that if \(\mathfrak{S}_1 \subset \mathfrak{S}_2\) and \(\mathfrak{S}_1 \neq \mathfrak{S}_2\) , then the uniformity of \(\mathfrak{S}_1\) -convergence is strictly coarser than the uniformity of \(\mathfrak{S}_2\) -convergence. In particular:
\end{enumerate}

\begin{enumerate}
\def\labelenumi{(\roman{enumi})}
\item
  If X is a non-compact Hausdorff topological space, the uniformity of compact convergence is strictly coarser than the uniformity of uniform convergence.
\item
  If X is a Hausdorff topological space which has infinite compact subsets (cf.~Chapter I, § 9, Exercise 4), then the uniformity of pointwise convergence is strictly coarser than the uniformity of compact convergence.
\end{enumerate}

\begin{enumerate}
\def\labelenumi{\arabic{enumi})}
\setcounter{enumi}{2}
\item
  Let X be a set, let S be a covering of X, let Y be a non-Hausdorff uniform space and let \(Y_{0}\) be the Hausdorff uniform space associated with Y (Chapter II, § 3, no. 8). Show that the Hausdorff uniform space associated with \(\mathcal{F}_{\mathfrak{S}}(\mathrm{X};\mathrm{Y})\) is isomorphic to \(\mathcal{F}_{\mathfrak{S}}(\mathrm{X};\mathrm{Y}_{0})\) .
\item
  Show that, on the set \(\mathcal{C}(\mathbf{R};\mathbf{R})\) of finite continuous real-valued functions defined on \(\mathbf{R}\) , the topology of uniform convergence with respect to the additive uniformity of R is not the same as the topology of uniform convergence with respect to the uniformity induced on R by the (unique) uniformity of the extended line \(\overline{R}\) (although the topologies induced by these two uniformities on R are the same).
\item
  Let X be a topological space. A set S of subsets of X is said to be saturated if it satisfies conditions \((\mathbf{F}_{\mathbf{I}}^{\prime})\) and \((\mathbf{F}_{\mathbf{II}}^{\prime})\) of no. 2 and if the closure of every set of S belongs to S.
\end{enumerate}

\begin{enumerate}
\def\labelenumi{\alph{enumi})}
\item
  Show that if X is normal and if S is saturated and covers X, then the set C(X; R) is dense in the space \(\tilde{\mathcal{C}}_{\mathfrak{S}}(X; \mathbf{R})\) (no. 6, Theorem 2, Corollary 2) with respect to the topology of S-convergence. In particular, C(X; R) is dense in \(F_{s}(X; \mathbf{R})\) . C(X; R) is closed in \(F_{s}(X; \mathbf{R})\) only if X is discrete.
\item
  Let X be a completely regular space (Chapter IX, § 1, no. 5) and let \(\mathfrak{S}_1, \mathfrak{S}_2\) be two saturated sets of subsets of X, such that \(\mathfrak{S}_1 \subset \mathfrak{S}_2\) and \(\mathfrak{S}_1 \neq \mathfrak{S}_2\) . Show that the topology of \(\mathfrak{S}_1\) -convergence on \(\mathcal{C}(\mathbf{X}; \mathbf{R})\) is strictly coarser than the topology of \(\mathfrak{S}_2\) -convergence.
\end{enumerate}

\begin{enumerate}
\def\labelenumi{\arabic{enumi})}
\setcounter{enumi}{5}
\tightlist
\item
  \begin{enumerate}
  \def\labelenumii{\alph{enumii})}
  \tightlist
  \item
    Let X be a topological space, Y a uniform space, and let \(\Phi\) be a filter on the set \(\mathcal{C}(X;Y)\) which converges pointwise to a function \(u_{0}\) . Then \(u_{0}\) is continuous at a point \(x_{0}\in X\) if and only if, for each entourage V of Y and each set \(M\in\Phi\) , there is a neighbourhood U of \(x_{0}\) and a mapping \(u\in M\) such that \((u_{0}(x),u(x))\in V\) for all \(x\in U\) .
  \end{enumerate}
\end{enumerate}

\begin{enumerate}
\def\labelenumi{\alph{enumi})}
\setcounter{enumi}{1}
\tightlist
\item
  Let X be a quasi-compact space, Y a uniform space and \(\Phi\) a filter on \(\mathcal{C}(X;Y)\) which converges pointwise to a function \(u_{0}\) . Then \(u_{0}\) is continuous on X if and only if, for each entourage V of Y and each set \(M\in\Phi\) , there exists a finite number of functions \(u_{i}\in\mathbf{M}(1\leqslant i\leqslant n)\) such that, for each \(x\in X\) , we have \((u_{0}(x),u_{i}(x))\in V\) for at least one index i.
\end{enumerate}

\begin{enumerate}
\def\labelenumi{\arabic{enumi})}
\setcounter{enumi}{6}
\tightlist
\item
  Let X, Y be two Hausdorff uniform spaces. For each continuous mapping \(f: \mathbf{X} \to \mathbf{Y}\) , let \(\mathbf{G}(f) \subset \mathbf{X} \times \mathbf{Y}\) be the graph of \(f\) ; \(\mathbf{G}(f)\) is closed in \(\mathbf{X} \times \mathbf{Y}\) (Chapter I, § 8, no. 1, Proposition 2, Corollary 2): hence \(f \to \mathbf{G}(f)\) is an injective mapping of \(\mathcal{C}(\mathbf{X}; \mathbf{Y})\) into the set \(\mathfrak{F}(\mathbf{X} \times \mathbf{Y})\) of non-empty closed subsets of \(\mathbf{X} \times \mathbf{Y}\) .
\end{enumerate}

\begin{enumerate}
\def\labelenumi{\alph{enumi})}
\item
  Show that the mapping \(f \to \mathbf{G}(f)\) of \(\mathcal{C}_u(\mathbf{X};\mathbf{Y})\) into \(\mathfrak{F}(\mathbf{X}\times \mathbf{Y})\) is uniformly continuous when \(\mathfrak{F}(\mathbf{X}\times \mathbf{Y})\) carries the uniformity defined in Chapter II, § 1, Exercise 5.
\item
  Let \(\Gamma\) be the image of \(\mathcal{C}(\mathbf{X};\mathbf{Y})\) in \(\mathfrak{F}(\mathbf{X}\times \mathbf{Y})\) under the mapping \(f\to \mathbf{G}(f)\) , and let \(\varphi :\Gamma \to \mathcal{C}_u(\mathbf{X};\mathbf{Y})\) be the inverse mapping. Show that if \(\mathbf{X}\) is compact, then \(\varphi\) is continuous on \(\Gamma\) (argue by contradiction).
\item
  Take both X and Y to be the compact interval {[}0, 1{]} of R. Show that \(\varphi\) is not uniformly continuous on \(\Gamma\) .
\end{enumerate}

\begin{enumerate}
\def\labelenumi{\arabic{enumi})}
\setcounter{enumi}{7}
\tightlist
\item
  Let X, Y be two metric spaces and let \((f_{n})\) be a sequence of Borel mappings of class \(\alpha\) of X into Y (Chapter IX, § 6, Exercise 16).
\end{enumerate}

\begin{enumerate}
\def\labelenumi{\alph{enumi})}
\item
  Suppose that the sequence \((f_{n})\) converges pointwise to a mapping f: X → Y. Show that f is of class \(\alpha + 1\) . {[}If U is an open subset of Y, note that \(\overline{f}^{1}(U) = \bigcup_{n \geqslant 1} \left( \bigcap_{k \geqslant 0} \overline{f}_{n+k}^{1}(U) \right)\) , and use Chapter IX, § 6, Exercise 4.{]}
\item
  Suppose that the sequence \((f_n)\) converges uniformly to \(f\) . Show that \(f\) is of class \(\alpha\) . {[}Let \(\mathbf{F}\) be a closed subset of \(\mathbf{Y}\) , and for each integer \(n > 0\) , let \(\mathbf{V}_n\) be the set of points of \(\mathbf{Y}\) whose distance from \(\mathbf{F}\) is \(\leqslant 1/n\) . Show that there is an increasing sequence \(n \to m(n)\) of integers such that \(\overline{f}^1(\mathbf{F}) = \bigcap_{n} \overline{f}_{m(n)}^1(\mathbf{V}_n)\) .{]}
\item
  Suppose that Y is of countable type. Show that, if \(f: X \to Y\) is a Borel function of class \(\alpha > 0\) , there exists a sequence of Borel functions \(g_{n}: X \to Y\) , of class \(<\alpha\) , such that \((g_{n})\) converges pointwise to f.~{[}Show that we may take the \(g_{n}\) to be functions which take only a finite number of values; use Exercises 4 h) and 16 c) of Chapter IX, § 6.{]}
\end{enumerate}

\begin{enumerate}
\def\labelenumi{\arabic{enumi})}
\setcounter{enumi}{8}
\item
  Let X be a Baire space, Y a metric space and \((f_{n})\) a sequence of continuous mappings of X into Y which converges pointwise to a function f.~A point \(x \in X\) is said to be a point of uniform convergence of the sequence \((f_{n})\) if, for each \(\varepsilon > 0\) , there exists a neighbourhood V of x and an integer p such that \(d(f_{m}(y), f_{n}(y)) \leqslant \varepsilon\) for all \(y \in V\) and all integers \(m \geqslant p, n \geqslant p, d\) being the metric on Y. Show that the complement S of the set of points of uniform convergence of the sequence \((f_{n})\) is meagre in X {[}cf.~Chapter IX, § 5, Exercise 22 a{]}. Give an example where S is dense in X. {[}Take X = Y = R, let \(n \to r_{n}\) be a bijection of N onto Q, and let \((g_{n})\) be a sequence of continuous mappings of R into {[}0, 1{]} which converges to the zero function for \(x \neq 0\) and to 1 at x = 0, the convergence being uniform in every open set of R which does not contain 0. Then consider the sequence of functions \(f_{n}(x) = \sum_{p=0}^{\infty} \alpha_{p} g_{n}(x - r_{p})\) , where the sequence \((\alpha_{p})\) tends to 0 in a suitable way.{]}
\item
  Show that, on the group \(\Gamma\) of homeomorphisms of the real line \(\mathbf{R}\) onto itself, the topology of pointwise convergence is the same as the topology of compact convergence (cf.~§ 3, Exercise 14).
\item
  Let X be a topological space and G a topological group. The set C(X; G) of continuous mappings of X into G is a subgroup of the group \(G^{X}\) . Let S be a set of subsets of X.
\end{enumerate}

\begin{enumerate}
\def\labelenumi{\alph{enumi})}
\item
  Suppose that for each \(A \in S\) , each neighbourhood V of e in G and each \(u \in C(X; G)\) , there exists a neighbourhood W of e in G such that \(sWs^{-1} \subset W\) for all \(s \in u(A)\) . Show that the topology of S-convergence is then compatible with the group structure of \(\mathcal{C}(X; G)\) , and that the right (resp. left) uniformity of the topological group \(\mathcal{C}_{\mathfrak{S}}(X; G)\) so defined is the same as the uniformity of S-convergence. Consider the case of pointwise convergence, the case of compact convergence when G is locally compact, and the case where G is abelian.
\item
  If G is the group \(\mathbf{SL}(2, \mathbf{R})\) endowed with the topology induced by that of \(R^{4}\) , show that the topology of uniform convergence is not compatible with the group structure of \(\mathcal{C}(\mathbf{R}; \mathbf{G})\) .
\end{enumerate}

\begin{enumerate}
\def\labelenumi{\arabic{enumi})}
\setcounter{enumi}{11}
\tightlist
\item
  Let X be a topological space and let A be a topological ring (Chapter III, § 6, no. 3). The set C(X; A) of continuous mappings of X into A is a subring of the ring A\^{}X. Let S be a set of subsets of X.
\end{enumerate}

\begin{enumerate}
\def\labelenumi{\alph{enumi})}
\item
  Suppose that, for each \(\mathbf{M} \in \mathfrak{S}\) and each \(u \in \mathcal{C}(\mathbf{X}; \mathbf{A})\) , the set \(u(\mathbf{M})\) is bounded (Chapter III, § 6, Exercise 12). Show that the topology of \(\mathfrak{S}\) -convergence is compatible with the ring structure of \(\mathcal{C}(\mathbf{X}; \mathbf{A})\) . Consider the case of pointwise convergence and the case of compact convergence.
\item
  Let X be a locally compact \(\sigma\) -compact space which is not compact (Chapter I, § 9, no. 9). Show that the topology of uniform convergence is not compatible with the ring structure of \(\mathcal{C}\) (X; R).
\end{enumerate}

